\begin{statement}[C.3.2; Proved: the equianharmonic point is neither an extremum]
Set $T_0=\sqrt3/2$, $\mu_0=(2-\rho)/3$, and
$p=\wp'_\rho(\mu_0)$. Then $p\ne0$, $p^2=-g_3(\rho)$, and, for
$\delta=\tau-\rho$ on the order-three root chart,
\begin{align*}
 \mu&=\frac{2-\tau}{3}
      +\frac{3p}{16\pi^2}\delta^2
      +\frac{3ip}{32\pi^2T_0}\delta^3+O(\delta^4),\\
 t&=\frac{ip^2}{\pi^3}\delta^3
        \left(1+\frac{3i}{2T_0}\delta+O(\delta^2)\right),\\
 \log N&=\tfrac14\log3+\frac7{192}\operatorname{Re}t
       -\frac{9\pi}{512T_0|p|^{2/3}}|t|^{4/3}
       +O(|t|^{5/3}).
\end{align*}
Thus
\[
 D_{\rm part}(t)=-\frac3{2\pi}\log48
       +\frac7{32\pi}\operatorname{Re}t
       -\frac{27}{256T_0|p|^{2/3}}|t|^{4/3}
       +O(|t|^{5/3}).
\]
The point $t=0$ is neither a local maximum nor a local minimum.
\end{statement}
\begin{proof}
We use K4 after checking its differentiation: put
$\ell=\partial_z\log\vartheta_1(\mu,\tau)$,
$a=\operatorname{Im}\mu/T$, $Z=\ell+2\pi ia$, and
$R=\mu'-a$. The theta heat equation and
$\wp=-\partial_z\ell-E$ give on the zero locus
\[
 \mu'=\frac{i\ell}{2\pi}-\frac{ig_2}{6\pi\wp'},\qquad
 R=\frac{iZ}{2\pi}-\frac{ig_2}{6\pi\wp'},\qquad
 \partial_\tau\log N=\frac{Z}{8\pi i}(Z+4\pi iR).
\]
These equations are equivalent to K4's factored formula and verify
its normalizations. At $\rho$ the unit endomorphism fixing the
nonzero order-three point forces the gradient of $U_\rho$ there to
vanish, so $Z=0$. Explicitly the unit $\rho-1$ fixes $\mu_0$ modulo
the lattice and has no fixed nonzero real covector. Hence
$a=\mu'=-1/3$ and $\ell=2\pi i/3$.

Along the zero locus, a further heat-equation differentiation gives
\[
 \ell'=\frac{i\wp'}{4\pi}+\frac{iEg_2}{6\pi\wp'}.
\]
The Ramanujan derivatives at $\rho$ are
\[
 g_2'=\frac{3ip^2}{\pi},\quad
 g_2''=\frac{5i}{T_0}g_2',\quad
 \frac{(\wp')'}{\wp'}=\frac{3i}{2T_0},\quad E=\frac\pi{T_0}.
\]
They give $\mu''=3p/(8\pi^2)$,
$\mu'''=9ip/(16\pi^2T_0)$ and hence the asserted $\mu$ expansion.
Expansion of $t=g_2^3/(g_2^3-27g_3^2)$ gives the displayed $t$
formula. Since $\partial_\tau a=R/(2iT)$, one has
\[
 Z=\frac{ip}{4\pi}\delta+O(|\delta|^2),\qquad
 R=\frac{3p}{8\pi^2}\delta+O(|\delta|^2).
\]
The derivative formula now gives the harmonic cubic
$\operatorname{Re}(7ip^2\delta^3/(192\pi^3))=(7/192)
\operatorname{Re}t+O(|\delta|^4)$.

For the next term, direct differentiation of the Gaussian gives
\[
 \partial_\tau\partial_{\bar\tau}\log N
       =-\frac\pi{2T}|R|^2
       =-\frac{9|p|^2}{128\pi^3T_0}|\delta|^2+O(|\delta|^3).
\]
After the harmonic cubic has been written as a function of $t$, the
order-three stabilizer permits only a radial term in degree four:
in the linearizing coordinate
$(\tau-\rho)/(\tau-\bar\rho)$ invariant monomials have exponents
whose difference is divisible by three. The coefficient of
$|\delta|^4$ is therefore $-9|p|^2/(512\pi^3T_0)$, since
$\partial_\delta\partial_{\bar\delta}|\delta|^4=4|\delta|^2$.
Converting $|\delta|^4$ to $|t|^{4/3}$ proves the expansion, including
the remainder. Its nonzero linear term in $\operatorname{Re}t$
has both signs, which proves the last assertion. In particular
superharmonicity at this root point does not make it a local maximum.
\end{proof}
